% Copyright (C) 2026 Yu Zhai. LPPL 1.3c or later.
\documentclass[11pt,a4paper]{article}
\usepackage[margin=18mm,headheight=14pt]{geometry}
\usepackage{unicode-math}
\setmathfont{Latin Modern Math}
\usepackage{physics}
\usepackage{finaldelim}
\usepackage{experimental/finaldelim-draw}
\FinalDelimDrawOff
\usepackage{booktabs,array,listings,fancyhdr,etoolbox}
\usepackage[colorlinks=true,linkcolor=blue!45!black,urlcolor=blue!45!black]{hyperref}
\lstset{basicstyle=\ttfamily\fontsize{8.5}{10.5}\selectfont,
 columns=fullflexible,keepspaces=true,breaklines=true,
 frame=single,rulecolor=\color{black!20},showstringspaces=false,
 aboveskip=4pt,belowskip=4pt}
\setlength\parindent{0pt}\setlength\parskip{6pt plus 1pt}
\emergencystretch=2em
\pagestyle{fancy}\fancyhf{}\fancyhead[L]{finaldelim}\fancyhead[R]{User manual}\fancyfoot[C]{\thepage}
% Listings use live example files; the classic listing adds display-only markup.
\newcommand\Example[2][]{\par\medskip\noindent
 \begin{minipage}[t]{.475\linewidth}\vspace{0pt}
 \textbf{LaTeX source}\par\smallskip
 \ifstrequal{#2}{classic-split}
  {\lstinputlisting[escapeinside={(*@}{@*)}]{manual-examples/classic-split-highlight.txt}}
  {\lstinputlisting{manual-examples/#2.tex}}
 \end{minipage}\hfill
 \begin{minipage}[t]{.495\linewidth}\vspace{0pt}
 \textbf{Result}\par\smallskip
 #1\input{manual-examples/#2.tex}
 \end{minipage}\par\medskip
 \FinalDelimSetup{niji=false}}
\ExplSyntaxOn
\cs_new_protected:Npn\ClassicMode
 {\cs_set_eq:NN\equation\fdNativeEquation
  \cs_set_eq:NN\endequation\fdNativeEndEquation
  \cs_set_eq:NN\left\__fd_tex_left
  \cs_set_eq:NN\right\__fd_tex_right
  \int_zero:N\l__fd_sizing_int}
\ExplSyntaxOff
\newcommand\StirlingNote{\footnote{Wikipedia,
 \href{https://en.wikipedia.org/wiki/Stirling\%27s_approximation\#Speed_of_convergence_and_error_estimates}{Stirling's approximation}, Stirling series.}}
\newcommand\NormalNote{\footnote{Wikipedia,
 \href{https://en.wikipedia.org/wiki/Normal_distribution\#Cumulative_distribution_function}{Normal distribution}, cumulative distribution function.}}
\begin{document}
{\LARGE\bfseries finaldelim}\par
{\large Delimiter sizing, nesting colors, and equation breaks}\par
Yu Zhai \quad\href{mailto:me@zhaiyusci.net}{me@zhaiyusci.net}\par
Version 0.1.0-dev \quad 5 November 2026

finaldelim adds three features to LuaLaTeX mathematics: distance-weighted
delimiter sizing, nesting colors named \texttt{niji}, and explicit line breaks
directly inside \texttt{equation}.

Compile with \textbf{LuaLaTeX} and load the package:
\begin{lstlisting}
\usepackage{unicode-math}
\setmathfont{Latin Modern Math}
\usepackage{finaldelim}
\end{lstlisting}
The new sizing policy is enabled by default. The manual uses Latin Modern Math;
finaldelim also works with other math fonts. Optional continuous-outline rendering
is described in Section~\ref{sec:outlines}. It loads \texttt{amsmath} and
\texttt{xcolor} automatically.

\section{Delimiter sizing}
Native \verb|\left| and \verb|\right| size a pair from its tallest enclosed
content. A distant fraction can therefore make both delimiters large beside
much smaller neighbours. finaldelim weights content by its distance from each
edge, balancing the pair with the surrounding expression.

The correction factor in Stirling's series mixes low terms with taller
fractions.\StirlingNote\ Here
$C(n)=n!/\bigl[\sqrt{2\pi n}(n/e)^n\bigr]$.
\par\textbf{Native sizing}
\Example{size-native}
\textbf{Default weighted sizing}
\Example{size-weighted}

\newpage
\subsection{Why distance, width, and height matter}
The model treats the contents as a horizontal profile. A measured box of height
$h_i$ and depth $d_i$ has radius
\[
 r_i=\max(h_i-a,\,d_i+a),
\]
where $a$ is the math axis height. Its radius applies over its actual width.
This gives a broad tall term more influence than a narrow one at the same
height. Spaces use the smallest measured radius in the group.

Let $r(x)$ be this profile over the content width $W$. The two delimiters see
opposite distance weights:
\[
 w_L(x)=e^{-x/\lambda},\qquad w_R(x)=e^{-(W-x)/\lambda}.
\]
At distance $\lambda$, a contribution has $1/e$ of its weight at the edge.
The left and right radius estimates are
\[
 R_L=\left(\frac{\int_0^W r(x)^q w_L(x)\,dx}
                    {\int_0^W w_L(x)\,dx}\right)^{1/q},\qquad
 R_R=\left(\frac{\int_0^W r(x)^q w_R(x)\,dx}
                    {\int_0^W w_R(x)\,dx}\right)^{1/q}.
\]
The larger estimate sets a shared target radius:
\[
 R=s\max(R_L,R_R),\qquad s=\texttt{scale}.
\]
The power $q$ gives tall content extra influence even when it occupies only a
small part of the width. Taking the maximum prevents a distant delimiter
from diluting the estimate at a delimiter close to tall content.
TeX then selects actual delimiters from the math font using this radius.

\subsection{Choosing settings}
\begin{center}
\begin{tabular}{@{}llp{94mm}@{}}\toprule
Option&Default&Effect\\\midrule
\texttt{lambda}&\texttt{20pt}&A shorter decay length makes nearby content more decisive.\\
\texttt{q}&\texttt{8}&A larger value gives tall content more influence.\\
\texttt{scale}&\texttt{0.9}&Multiplies the target radius about the math axis.\\\bottomrule
\end{tabular}
\end{center}
The default decay length is a physical distance: 20pt at every font size.
\texttt{lambda} and \texttt{scale} must be positive; \texttt{q} must be nonnegative.
The special value \texttt{q=0} selects the maximum-radius limit.

\begin{lstlisting}
\usepackage[lambda=20pt,q=8,scale=0.9]{finaldelim}
% Change settings for subsequent formulas:
\FinalDelimSetup{scale=1} % use the full estimated radius
% Return to native sizing; colors and line breaks remain available:
\FinalDelimSetup{sizing=native}
\FinalDelimSetup{sizing=weighted,lambda=20pt}
\end{lstlisting}
For a local setting, write:\par
\verb|\begin{equation}[lambda=30pt]|

The default scale of 0.9 favours a compact visual frame rather than requiring
the delimiters to enclose the full height of the contents.
The scale applies to weighted and endpoint sizing. It changes the sizing
target, not the glyph outline; font size steps can leave a glyph unchanged.
Native sizing and explicitly sized delimiters are unaffected.

\newpage
\subsection{Continuous delimiter outlines}\label{sec:outlines}
The sizing model chooses how large a delimiter should be. Continuous-outline
rendering gives it intermediate sizes between the font's prebuilt glyphs and
fills assembled delimiters as a single outline, avoiding repeated painting at
component joins. It uses the current math font's own shapes.

Enable this optional experimental feature after loading finaldelim:
\begin{lstlisting}
\usepackage{finaldelim}
\usepackage{experimental/finaldelim-draw}
\end{lstlisting}
Keep \texttt{finaldelim-draw.sty}, \texttt{finaldelim-draw.lua} and
\texttt{finaldelim-outline-fonts.lua} in an \texttt{experimental/} subdirectory
beside the document. The renderer loads font outlines automatically.

Then write ordinary \verb|\left|, \verb|\middle| and \verb|\right| commands.
Niji and line breaks work with the drawn outlines too.
Superscripts and subscripts after a visible closing delimiter follow its final
size; the contents retain their full height and depth for surrounding layout.

This four-by-two matrix illustrates the Frobenius norm.\footnote{Wikipedia,
\href{https://en.wikipedia.org/wiki/Matrix_norm\#Frobenius_norm}{Matrix norm},
Frobenius norm; shown here for a four-by-two matrix.}
Both results below use the same default sizing settings.
\par\textbf{Native font glyphs}
\Example[\FinalDelimDrawOff]{outline-norm}
\textbf{Continuous outlines}
\Example[\FinalDelimDrawOn]{outline-norm}

\verb|\FinalDelimDrawOff| returns to native glyph rendering;
\verb|\FinalDelimDrawOn| enables drawing again. These commands respect TeX groups
and do not change the sizing settings.

Between prebuilt sizes, the renderer interpolates the font's contours.
For larger assembled symbols it keeps the font's component placement but paints
one combined path. Slashes, backslashes, single angles and double angles can
extend beyond the largest prebuilt glyph: their width follows the recent size
trend, while stroke thickness and terminal details are preserved. Double angles
retain their inter-stroke spacing. Other symbols without larger variants or
assembly parts retain their native size limit. Transformed fonts and explicit
variable-font instances currently retain native rendering. Drawn delimiters are
PDF paths; text extraction and accessibility tagging for these paths are not yet supported.

\newpage
\subsection{One or more middle delimiters}
Use \verb|\middle| for a delimiter inside a \verb|\left|/\verb|\right| group.
Every middle delimiter participates in the size calculation. A matrix element
in bra--ket notation uses two of them.\footnote{Wikipedia,
\href{https://en.wikipedia.org/wiki/Bra\%E2\%80\%93ket_notation\#Linear_operators_acting_on_bras}{Bra--ket notation}, linear operators acting on bras.}
\Example{middle-pair}
The outer delimiters and both bars share one sizing target. The math font
supplies the actual glyphs. With \texttt{niji}, the middle delimiters receive
the same color as their surrounding pair.

\subsection*{Each delimiter sees its own neighbourhood}
Let $s_j$ be the horizontal position of delimiter $j$, including both ends
and every middle delimiter. Its weight and radius estimate are
\[
 w_j(x)=e^{-|x-s_j|/\lambda},\qquad
 R_j=\left(\frac{\int_0^W r(x)^q w_j(x)\,dx}
                  {\int_0^W w_j(x)\,dx}\right)^{1/q}.
\]
For $N$ delimiters, their shared radius is
\[
 R=s\max_{1\le j\le N}R_j.
\]
A middle delimiter receives contributions from both sides. Moving it closer
to tall content can increase its estimate and therefore the shared size.
With only the two outer delimiters, this is the earlier left/right model.

\subsection*{Across line breaks}
Middle delimiters stay at their source positions when the formula is split
into rows. Under \texttt{spanning=endpoints}, the size envelope includes every
row carrying one of the group's delimiters: the first row, the last row,
and any row containing a middle delimiter. Under \texttt{spanning=full}, all
delimiter positions use the distance-weighted model over the continuous contents.
Nested groups are sized independently.

\newpage
\section{Niji: colors for nested delimiters}
Enable \texttt{niji} to give matching delimiters the same nesting color:
\begin{lstlisting}
\usepackage[niji]{finaldelim}
\end{lstlisting}
\texttt{rainbow} is an alias. Colors apply to explicit \verb|\left|/\verb|\right|
pairs; a \verb|\middle| delimiter follows its pair. Ordinary parentheses and
the mathematical contents keep their usual color.

The normal cumulative distribution function has a bracketed expression with
an inner argument of the error function.\NormalNote
\par\textbf{Without niji}
\Example{niji-plain}
\textbf{With niji}
\Example{niji-display}

\subsection{Inline mathematics}
The same setting works in ordinary inline math.
\Example{niji-inline}

\subsection{Settings and palettes}
\begin{lstlisting}
\FinalDelimSetup{niji=false}             % subsequent mathematics
\FinalDelimSetup{niji,colors={blue,red,teal}}
\end{lstlisting}
Use an environment option, such as \verb|\begin{equation}[niji=false]|,
to override the setting for one display. The palette repeats at deeper nesting
levels. Sizing and coloring can be selected independently.

\newpage
\section{Line breaks directly in equation}
Place \verb|\\| at a break and \verb|&| at the alignment point.
Keep the original delimiter pair: \verb|\left| can open on one row and
\verb|\right| can close on another.

\subsection{The same long formula, old and new}
The divergence between two multivariate normal distributions contains a trace,
a quadratic form, and a logarithmic determinant.\footnote{Wikipedia,
\href{https://en.wikipedia.org/wiki/Kullback\%E2\%80\%93Leibler_divergence\#Multivariate_normal_distributions}{Kullback--Leibler divergence}, multivariate normal distributions.}
Keeping the full notation makes this a useful case for several rows: the
left-hand side introduces the quantity, and each subsequent row carries a term.

\textbf{Plain amsmath: split and invisible delimiters}\par
The traditional form uses an inner \texttt{split}. The opening bracket is
paired with \verb|\right.|; the closing bracket is paired with \verb|\left.|.
The intermediate row needs no delimiter commands. A \verb|\vphantom| copies
the height and depth of the final logarithmic fraction into the opening row,
so both brackets receive the same size.
\textcolor{red!70!black}{\textbf{Red bold}} marks the extra code.
\Example[\ClassicMode]{classic-split}

\textbf{finaldelim: direct breaks, with niji included}\par
Keep the same break points, remove the extra wrapper, invisible delimiters, and height phantom,
and keep a single
\verb|\left|/\verb|\right| pair across the rows.
\Example{direct-break}
Niji is included as an optional companion feature: adding \texttt{[niji]}
colors matching delimiters by nesting level, including pairs that cross rows.
The example also uses \verb|\left|/\verb|\right| around the trace argument
so its inner pair receives the next color. Both forms keep one equation number.

\newpage
\subsection{Break a long expansion into readable groups}
The cubic trinomial expansion has ten terms.\footnote{Wikipedia,
\href{https://en.wikipedia.org/wiki/Multinomial_theorem\#Example}{Multinomial theorem}, cubic trinomial example.}
Breaking before groups of mixed terms makes it easier to scan while preserving
one equation number. The source contains only \texttt{equation}, alignment
markers, and the chosen breaks.
\Example{multinomial-break}

\subsection{Break at a change in the calculation}
The normal-distribution log-likelihood starts with a sum and then expands into
logarithmic terms and a sum of squared residuals.\footnote{Wikipedia,
\href{https://en.wikipedia.org/wiki/Normal_distribution\#Estimation_of_parameters}{Normal distribution}, estimation of parameters.}
A new row separates the expansion step; another keeps the residual sum together.
\Example{loglike-break}

\subsection{Choose the spanning size}
\texttt{spanning=endpoints}, the default, uses the maximum height and depth of
the first and last rows inside a spanning pair, plus any rows carrying its
\verb|\middle| delimiters. \texttt{spanning=full} instead
applies the weighted model to its entire continuous contents, including the
middle rows. Select either policy on the display:
\begin{lstlisting}
\begin{equation}[spanning=full,niji]
 ...
\end{equation}
\end{lstlisting}
For pairs contained on one row, weighted sizing applies in either mode.
\verb|\\[3pt]| adds space after a row; \verb|\jot| controls the usual extra
row spacing. Additional \verb|&| markers create additional alignment pairs.

\newpage
\subsection{A number for each row}
Change \texttt{equation} to \texttt{align} when rows need separate numbers.
The maximum-likelihood estimates of a normal distribution provide a short
example.\footnote{Wikipedia,
\href{https://en.wikipedia.org/wiki/Normal_distribution\#Estimation_of_parameters}{Normal distribution}, estimation of parameters.}
\Example{numbered}
The labels refer separately to the mean, \eqref{eq:mean}, and the variance,
\eqref{eq:variance}. Use \verb|\notag| on a row to suppress its number.
\texttt{align*} suppresses all automatic numbers.

\subsection{Environment guide}
\begin{center}
\begin{tabular}{@{}lp{100mm}@{}}\toprule
Environment&Use\\\midrule
\texttt{equation}&Direct breaks and alignment, one number for the whole formula.\\
\texttt{equation*}, \verb|\[...\]|&The same layout without an automatic number.\\
\texttt{align}, \texttt{align*}&With or without per-row numbering.\\
\texttt{alignat}, \texttt{alignat*}&An explicit column-pair count, with user-controlled spacing.\\
\texttt{matchedaligned}&An embedded alignment, with optional \texttt{[t]}, \texttt{[c]}, or \texttt{[b]} positioning.\\\bottomrule
\end{tabular}
\end{center}
The explicit names \texttt{matchedsplit}, \texttt{matchedalign} and
\texttt{matchedalignat} are also available. Multiple alignment pairs are
supported; the author determines their mathematical interpretation.

\subsection{Layout controls}
finaldelim lays out the mathematics continuously and then separates the actual
contents at the row and alignment markers. Fractions, scripts and embedded
matrices stay together. You choose where to break.

Use \texttt{debug=boxes}, \texttt{debug=alignment}, or \texttt{debug=all} as a
package or display option to inspect boxes and alignment points. Box outlines
use \texttt{lua-visual-debug} and its \verb|\lvdset| interface.

\newpage
\subsection{Using physics}
Load \texttt{physics} and \texttt{finaldelim} in either order. The integration
is automatic; no extra option is needed. Automatic delimiters created by
\verb|\qty|, \verb|\abs|, \verb|\norm| and bra--ket commands use the same sizing
and niji colours as explicit \verb|\left| and \verb|\right| pairs.

The Gaussian relative-entropy example can also be written with physics
commands.\footnote{Wikipedia,
\href{https://en.wikipedia.org/wiki/Multivariate_normal_distribution\#Kullback\%E2\%80\%93Leibler_divergence}{Multivariate normal distribution}, Kullback--Leibler divergence.}
Here the breaks and alignment points are inside the argument of \verb|\bqty|:
\Example{physics-break}

Direct breaks are recognized in the automatic forms of \verb|\qty|,
\verb|\quantity|, \verb|\pqty|, \verb|\bqty|, \verb|\Bqty|, \verb|\vqty|,
\verb|\abs|, \verb|\absolutevalue| and \verb|\norm|. Matrix arguments of
\verb|\mqty|, \verb|\matrixquantity|, \verb|\smqty| and
\verb|\smallmatrixquantity| retain their own rows and columns. For example,
\verb|\mqty(a&b\\c&d)| remains a two-by-two matrix inside an equation.

Starred forms keep physics' fixed-size behaviour. Their arguments, fractions,
subscripts and superscripts remain indivisible. Other macros are not expanded
by the break scanner; use explicit delimiters for breaks hidden inside an
unsupported macro's argument. The normalization pass does not evaluate argument contents; they are left
for mathematical typesetting.

\subsection*{Installation}
Install these six files together beside your document or in a local TEXMF tree:
\path{finaldelim.sty}, \path{finaldelim-luatex.sty},
\path{finaldelim-interface.sty}, \path{finaldelim-physics.sty},
\path{finaldelim.lua}, and
\path{finaldelim-size.lua}. This implementation targets LuaLaTeX; tagged
mathematical structure for the physical line-breaking backend is not yet supported.

Maintained by Yu Zhai. LPPL 1.3c or later. Formula sources are given in footnotes;
notation and line breaks are adapted for these examples. Sources consulted
4 November 2026.
\end{document}
